\documentclass[12pt,answers]{exam}


\usepackage[margin=0.8in,footskip=0.2in]{geometry}

\usepackage{etex}
\usepackage{amssymb,amsmath,multicol} %<-- InWorksheetExam1 i also have fancyhdr,

\usepackage[metapost]{mfpic}
\usepackage[pdftex]{graphicx}

\usepackage{pst-plot}
\usepackage{pgfplots}
\pgfplotsset{compat=1.9}

\usepackage{tikz}
\usepackage{tkz-2d}
\usepackage{tkz-base}
\usetikzlibrary{calc}

\usepackage{tabularx, booktabs}
\newcolumntype{Y}{>{\centering\arraybackslash}X}

\usepackage[inline]{enumitem}
\usepackage{refcount}%<-- non in WorksheetExam1

\usepackage{pstricks-add,pst-eucl}
\usepackage{caption}

\def\f{x+1} \def\g{-x/3+2}  \def\h{-x+3}

\newcommand{\vasymptote}[2][]{
	\draw [densely dashed,#1] ({rel axis cs:0,0} -| {axis cs:#2,0}) -- ({rel axis cs:0,1} -| {axis cs:#2,0});
}

\newcommand*\circled[1]{\tikz[baseline=(char.base)]{%
		\node[shape=circle,draw,inner sep=1pt] (char) {#1};}}
\renewcommand\choicelabel{\circled{\thechoice}}


\let\originalleft\left
\let\originalright\right
\renewcommand{\left}{\mathopen{}\mathclose\bgroup\originalleft}
\renewcommand{\right}{\aftergroup\egroup\originalright}

\addpoints
%\printanswers
\noprintanswers

\usepackage{geometry}
\geometry{
	%a4paper,
	%total={170mm,257mm},
	%left=20mm,
	%top=20mm,
	text={.95\paperwidth,.95\paperheight}, ratio=1:3,%includefoot
}

\opengraphsfile{Q10a_Spring_16}

\begin{document}
%\extrawidth{-0.3in}
\pagestyle{headandfoot}

%\setlength{\hoffset}{-.25in}

%\extraheadheight{-.3in}
\runningheadrule
\firstpageheader{\bfseries {MATH1-UC 1171}}{ \bfseries {Quiz 10 }}{\bfseries {4/19/2016}} 



\firstpagefooter{Bonus question on back of this page} %%&&CHANGED
                {}
                {Points earned: \hbox to 0.5in{\hrulefill}
                 out of  \pointsonpage{\thepage} points}
                 
						

%\vspace*{0.1cm}
\hbox to \textwidth { \scshape {Name:} \enspace\hrulefill}
\vspace{0.1in}




\pointpoints{point}{points}

\begin{questions}


\addpoints

\question The population of a certain city was 110,000 in 2006, and the observed doubling time for the population is 17 years.
\begin{parts}
\part[2] Find an exponential model 
$\displaystyle n(t) = n_02^{\frac{t}{a}}$
 for the population $t$ years after 2006. (Show your work step-by-step.)
\fillwithdottedlines{2.5cm}
\part[2] Another was to write the formula for the population $n(t)$ is by using the exponential function with base $e$. The formula is $\displaystyle n(t)=n_0 e^{rt}$. Find $r$. (Show your work step-by-step and do not simplify your answer.)
\fillwithdottedlines{2.5cm}
\end{parts}
\question[2] The half-life of radium-226 is 1600 years. Suppose we have a 28-mg sample.
Find a function 
$\displaystyle m(t) = m_02^{-\frac{t}{h}}$
 that models the mass remaining after $t$ years. (Show your work step-by-step.)
 \fillwithdottedlines{2.5cm}
 \question[2] A patient takes a medicine that has a half-life of 2 hours (that is, it takes 2 hours for concentration of the medicine in the patient's system to drop by 50\%.) An overdose occurs when the patient has a  concentration of 190mg/l (milligrams per liters of blood) of medicine in her system. Every time the patient takes a pill containing this medicine, the concentration of the medicine in the patient's system increases by 100mg/l.

At 1 pm, the patient has just taken the first dosage of the medicine, so she has a concentration of 100mg/l of medicine in her system. She takes the same amount of medicine every two hours. When will she overdose? Explain your reasoning and show your work step-by-step. 
\fillwithdottedlines{4cm}
\newpage
\thispagestyle{empty}
 \bonusquestion[2] Last week we created a new number line, where the tick mark 1 represented the exponent of $\displaystyle 10^1$, the tick mark 2 represented the exponent of $\displaystyle 10^2$, and so on. In this number line, we had no difficulty visualizing both 1,000,000 (the tick mark 6) and 1,000,000 (the tick mark 9.) We noted that this is a \emph{logarithmic} scale, because the exponents we plotted were logarithms ($\displaystyle 1=\log(10)$, $2=\log\left( 10^2\right)$, etc.) A way to label the logarithmic scale is with the values themselves instead of the exponents, as shown below. 
	
	\begin{tikzpicture}
	\draw (0.5,0) -- (11.5,0);
	
	\draw[xshift=2 cm] (0pt,2pt) -- (0pt,-2pt) node[below,fill=white] {0.1};
	\draw[xshift=4 cm] (0pt,2pt) -- (0pt,-2pt) node[below,fill=white] {1};
	\draw[xshift=6 cm] (0pt,2pt) -- (0pt,-2pt) node[below,fill=white] {10};
	\draw[xshift=8 cm] (0pt,2pt) -- (0pt,-2pt) node[below,fill=white] {100};
	\draw[xshift=10 cm] (0pt,2pt) -- (0pt,-2pt) node[below,fill=white] {1000};
	
	%\node at (6.5,-1.0) {Log scale};
	\end{tikzpicture}
	
	In this number line, each tick mark is 10 times as large as the tick mark immediately preceding it. Explain how you would plot the number 7 accurately on the number line shown above. Show all your work step-by-step. (Note: you will have to set up and solve an equation to receive credit.)
\fillwithdottedlines{6cm}	


	
%%%%%%%%%%%%%%%%%%%%%%

\end{questions}

\end{document}                 